\chapter{Two-Dimensional Defect Extraction and Material-Transfer Program}
\index{impurity operator!two-dimensional program|textbf}
\index{warmup!two-dimensional defect extraction|textbf}

\label{app:two-dimensional-defect-extraction}

\chapterstatus{active controlled synthetic defect-extraction program and
proposed material-transfer program. The periodic two-dimensional parent,
spin-space contracts, and one-dimensional defect-extraction chain are completed
prerequisites. The version-1 design is human-accepted, and the bounded synthetic
Stage~A null-and-folding calculation is human-accepted and independently
reproduced. The bounded synthetic Stage~B scalar-onsite multi-route calculation
is complete, independently reproduced, and human-accepted at its exact bounded
evidence boundary. The accepted-parent Stage~C software contract is implemented,
independently verified only with authored synthetic fixtures, and human-accepted
at that exact execution-free boundary; it has not read an accepted parent. No
accepted-parent directional or nonlocal two-dimensional defect result, material impurity
operator, continuum crossover, scientific validation, or uncertainty-
quantification result is claimed.}

\section{Purpose and prerequisite evidence}

This appendix defines the next controlled extension of the impurity-extraction
program. Appendix~\ref{app:two-dimensional-reduction} supplies the accepted
separable and coupled periodic parents. Appendix~\ref{app:impurity-model-classes}
supplies the accepted spin-space, alignment, extraction-route, finite-size,
model-class, analytical-oracle, and continuum-refinement contracts. The
controlled synthetic task is active. The accepted version-1 design and Stage~A
evidence are retained under
\texttt{calculations/research-monograph/impurity-defect-2d/}; they fix staged
synthetic inputs, represented-space maps, metrics, adverse controls, tolerances,
stopping rules, independent-verification routes, authority, and resource
bounds. The exact Stage~A execution record binds the design, runner, scalar
parent, canonical repository, output, resource envelope, and durable human
decision; its separate acceptance record applies only to this bounded synthetic
evidence. Stage~B initially exposed an incompatibility among the chosen generic
twist, componentwise $[0,1)$ twist reduction, and bare-permutation
full-Hamiltonian covariance. That single-route convention is retained as
negative design evidence. The adopted execution-free replacement constructs
independent centred uniform-link and reduced seam matrices, relates them through
one explicit gauge bridge, and executes A-then-B and B-then-A in fresh processes
to expose hidden schedule state. Its authored toy retains all 2,304 blind
candidates; clean schedules agree, an intentionally shared-cache mutation is
discriminating, and an independent verifier reconstructs the toy result. A
separately authorized calculation then applied the frozen two-route protocol to
the accepted isotropic scalar parent. It retained 72 known-map route
evaluations, 36 bridge records, and 2,304 blind candidate rows; runner criteria
and independent Fourier/seam reconstruction pass. This exact calculated
synthetic numerical-verification evidence is human-accepted, but it is not
material validation. Accepted-parent Stages~C--E remain defined only at program
level. The program adds
only genuinely two-dimensional defect structure; material transfer
remains proposed and requires separate authorization and retained calculated
evidence.

\section{Completed Stage~A null and folding control}

The human-accepted Stage~A calculation used the separable scalar parent and the
frozen Cartesian product of $6\times6$ and $8\times8$ cells with four boundary twists. All eight
cases passed the declared finite-matrix criteria. The maximum folding-unitarity
defect was $1.280\times10^{-15}$; the maximum folded off-block and primitive-block
defects were both $2.584\times10^{-16}$; the maximum eigenvalue defect was
$2.221\times10^{-16}E_G$; and the maximum recovered-null Frobenius defect was
$2.179\times10^{-16}E_G$. These are calculated synthetic floating-point
residuals under the frozen representation, not physical uncertainties.

An independently implemented Kronecker seam/Fourier route reconstructed every
retained metric and reported a maximum retained-versus-reconstructed difference
of $3.473\times10^{-16}$. The analytical seam controls separately reproduced
$1$, $\exp(+2\pi i\phi)$, and $\exp(-2\pi i\phi)$ for noncrossing, positive-crossing,
and negative-crossing hops at $\phi=0.37$ turns. Incorrect geometry, incorrect
boundary phase, lost site correspondence, and unknown energy reference each
produced the frozen structured stop before residual formation. This evidence
verifies the Stage~A folding, seam, null, and compatibility software contracts;
it did not itself activate Stage~B.

\section{Completed Stage~B scalar-onsite and symmetry control}

The Stage~B protocol fixes the integer $D_4$ matrices, the eight-site orbit of
an off-axis scalar onsite plant, known-map extraction, transformed-twist
full-Hamiltonian covariance, and energy-reference adverse controls. Its adopted
execution-free construction reconciles the finite-matrix twist gauges through
independent centred uniform-link and reduced seam routes plus one explicit
bridge. Its adversarial design analysis also records an identifiability limit:
an off-diagonal host-only blind objective cannot determine an absolute
translation in a homogeneous periodic parent. The blind route must therefore
retain the complete symmetry-related ambiguity set and stop before extraction
rather than select the planted site. The calculation reproduces that prediction:
the Gamma case retains 512 tied maps per route and schedule, while the generic
twist retains 64 tied translations. Every case stops with
\texttt{DEFECT\_2D.SITE\_MAP\_UNRESOLVED}, with null selected map and null extracted
operator.

Across both routes and schedules, the largest known-map recovery
maximum-entry defect is $2.238\times10^{-16}E_G$, the largest recovery Frobenius
defect is $3.079\times10^{-16}E_G$, and the largest attacked bridge Frobenius
residual is $3.537\times10^{-16}E_G$. Known-case and bridge schedule differences
are exactly zero in the retained comparisons. The independent verifier reports
a maximum reconstructed difference of $7.845\times10^{-13}$, below the frozen
$10^{-10}$ criterion. Omitting the known $0.137E_G$ shift produces a
$1.096E_G$ Frobenius residual, while the fixed-twist adverse controls produce
maximum-entry residuals $9.979\times10^{-3}E_G$ and
$4.065\times10^{-2}E_G$ in Routes A and B. These adverse outcomes remain
separate from the passing correctly transformed cases.

This result verifies the bounded finite synthetic scalar-onsite, $D_4$, gauge-
bridge, schedule, ambiguity, and adverse-control contracts only. It does not
support a rerun, material transfer, or a directional accepted-parent result.

A subsequent execution-free Stage~C software package uses authored bond
coefficients without reading the accepted periodic parent. It distinguishes a
directional nearest-neighbour plant from an isotropic class and a diagonal
nonlocal plant from all shorter-range classes. Two independently constructed
gauge routes, two fresh spawned schedule processes, all eight oriented $D_4$
operations at Gamma, five ordered model classes, locality shells, and four
adverse controls are retained by the toy wire contract. Maintained tests and an
independent implementation reproduce the toy behavior with maximum difference
$1.388\times10^{-17}E_G$. This exact execution-free package is human-accepted
at its bounded software-verification boundary. It is not an accepted-parent
Stage~C calculation or numerical evidence for a material.

The human-adopted accepted-parent Stage~C design binds the accepted isotropic
and anisotropic scalar parents. It separates isotropic $D_4$ covariance,
anisotropic $D_2$ covariance, and an axis swap to the distinct
$(\lambda_x,\lambda_y)=(0.7,0.3)$ parent. A separately authorized execution-free
implementation uses only an authored 61-hopping isotropic fixture and an authored
$15\times15$ anisotropic energy grid. Across two fresh processes it retains the
frozen 208 route evaluations, 104 gauge bridges, 1,040 ordered model fits, and
104 cross-schedule comparisons. The independent inverse-Fourier and QR verifier
reconstructs all records with maximum scalar difference
$1.922\times10^{-15}E_G$. Invalid anisotropic $D_4$ invariance and omission of
the separately swapped parent each produce $6.0\times10^{-3}E_G$, above their
predeclared $10^{-3}E_G$ adverse floors. A synthetic run completed in 1.56
seconds with 99,041,280 bytes maximum resident set size and a 2,205,921-byte
result, within the adopted envelope. Parent hopping reconstruction and
truncation, represented-gauge error, alignment recovery, and defect-model
residuals remain separate. This implementation has no accepted-parent input
adapter, has read no accepted parent, awaits separate human acceptance, and
supplies no calculated Stage~C evidence or Stage~D authority.

\section{Matched pristine and defect supercells}

Use the separable and coupled two-dimensional parents developed in Appendix~\ref{app:two-dimensional-reduction}. Let
the supercell contain $N_x\times N_y$ primitive cells with lattice vectors
\begin{equation}
    \mathbf L_1=N_x\mathbf a_1,
    \qquad
    \mathbf L_2=N_y\mathbf a_2.
\end{equation}
The matched pristine and defect fibers are
\begin{align}
    \mathbf H_{b,\mathbf L}(\mathbf K),
    \qquad
    \mathbf H_{d,\mathbf L}(\mathbf K),
\end{align}
with $\mathbf K$ in the reduced supercell Brillouin zone.

The pristine control must reproduce two-dimensional folding from the primitive
mesh into each supercell fiber. In the separable limit, the supercell
construction should also reproduce products and sums of the verified
one-dimensional results. Degenerate folded clusters must be compared through
their complete projectors rather than through arbitrary eigenvector labels.

The first planted defect should be a scalar onsite or localized multiplication
term at a declared supercell site. Accepted-parent stages should add:
\begin{enumerate}
    \item direction-dependent hopping changes;
    \item a finite-range nonlocal defect;
    \item anisotropic spatial profiles; and
    \item a composite-orbital defect in the coupled two-dimensional parent.
\end{enumerate}
Independent point-group-related placements should produce aligned operators
related by the corresponding lattice symmetry. A failure of this covariance
should be attributed among site mapping, supercell origin, gauge, degeneracy
handling, and boundary phases before the defect model is interpreted.

The finite-size study must vary both area and shape. Square sequences probe
increasing image separation under preserved symmetry, while rectangular
sequences test anisotropic image interactions and alignment conventions. One
apparently converged square supercell is not sufficient evidence for an
isolated two-dimensional defect.

\section{Energy-reference, gauge, ordering, and geometry stress tests}

For a known aligned defect, construct a raw candidate representation
\begin{equation}
    \widetilde{\mathbf H}_{d}
    =
    \mathbf G
    \left(
        \mathbf H_b+\mathbf V_{\mathrm{known}}
    \right)
    \mathbf G^\dagger
    +
    c\mathbf I,
\end{equation}
where $\mathbf G$ is a known admissible coordinate transformation and $c$ is a
known scalar shift. The transformation sequence should include:
\begin{enumerate}
    \item orbital phases;
    \item orbital and site permutations;
    \item unitary rotations within degenerate subspaces;
    \item primitive or supercell translations;
    \item admissible point-group transformations;
    \item changes of localized gauge; and
    \item a scalar energy-reference offset.
\end{enumerate}
After applying the inverse declared alignment and energy-reference procedure,
the extracted operator should equal $\mathbf V_{\mathrm{known}}$ within the
stated representation error.

The tests should also contain deliberately incompatible cases: unequal retained
rank, mismatched spin space, incorrect supercell shape, lost site
correspondence, or large principal angles between retained subspaces. These
cases should stop the comparison. They should not be forced into a unitary fit
merely because two final matrices can be made equal in dimension.

The energy-reference test is particularly important. Without correction,
$c\mathbf I$ appears as a spatially extended false impurity term. The alignment
procedure must remove only the declared reference offset and must not subtract
a genuine localized onsite component.

\section{Supercell-size and shape convergence}

Let $\Delta\mathbf H_L$ denote an aligned impurity operator for supercell size
or shape label $L$. Define a projector $\mathbf P_r$ onto a declared region
within distance $r$ of the defect and
\begin{equation}
    \mathbf Q_r=\mathbf I-\mathbf P_r.
\end{equation}
Complementary diagnostics include
\begin{align}
    \varepsilon_{\mathrm{core}}(r;L,L')
    &=
    \left\lVert
        \mathbf P_r
        \left(
            \Delta\mathbf H_L-
            \Delta\mathbf H_{L'}
        \right)
        \mathbf P_r
    \right\rVert_{\mathrm F},
    \\
    \varepsilon_{\mathrm{ext}}(r;L)
    &=
    \left\lVert
        \mathbf Q_r
        \Delta\mathbf H_L
        \mathbf Q_r
    \right\rVert_{\mathrm F},
    \\
    \varepsilon_{\mathrm{cross}}(r;L)
    &=
    \left(
        \left\lVert
            \mathbf P_r
            \Delta\mathbf H_L
            \mathbf Q_r
        \right\rVert_{\mathrm F}^2
        +
        \left\lVert
            \mathbf Q_r
            \Delta\mathbf H_L
            \mathbf P_r
        \right\rVert_{\mathrm F}^2
    \right)^{1/2}.
\end{align}
The comparison between different supercells requires an explicit embedding or
restriction to common site and orbital coordinates. The notation above assumes
that this identification has already been performed.

A global Frobenius norm can scale with represented volume or the number of
repeated image contributions. It must not replace local, exterior, and
cross-coupling diagnostics. At finite supercell size, defect-derived energies
must be tracked as functions of supercell momentum. Their bandwidth and
$\mathbf K$ dependence diagnose image coupling and should not be collapsed into
a single nominal bound-state energy. Wavefunction localization and
symmetry-resolved blocks should be tracked independently over the same size and
shape sequence.

\section{Recovery of known model classes}

Let the planted impurity belong exactly to a declared class
$\mathcal M^{(m_0)}$. Fit the ordered hierarchy
\begin{equation}
    \mathcal M^{(0)},
    \mathcal M^{(1)},
    \ldots
\end{equation}
to the extracted operator. The experiment should verify that:
\begin{enumerate}
    \item a scalar onsite class recovers a scalar onsite defect;
    \item an orbital-dependent class is required for orbital splitting;
    \item an onsite class fails for a planted hopping correction;
    \item a local class leaves a nonzero residual for a planted nonlocal
          operator;
    \item a spin-independent class fails for spin splitting or mixing; and
    \item the first class containing the planted structure can recover it
          within the independently established extraction error.
\end{enumerate}
The failure of a simpler class is a desired outcome when the planted defect lies
outside it. This test determines whether the minimum operator separation acts
as an expressiveness diagnostic rather than merely as an optimization score.

The model hierarchy, parameter bounds, weighting, and stopping rule must be
frozen before the outcomes are inspected. An unnecessarily richer class may
fit the planted operator but does not demonstrate that the selection rule can
identify the least adequate class.

\section{Smoothness-to-continuum crossover}

Introduce a defect profile with tunable width $\sigma$. One fixed-peak family
is
\begin{equation}
    V_{\sigma}^{(\mathrm{peak})}(\mathbf r)
    =
    V_c
    \exp\left(
        -\frac{\lVert\mathbf r\rVert^2}{2\sigma^2}
    \right).
\end{equation}
Its integrated strength changes with $\sigma$. A separate fixed-integrated
family in spatial dimension $D$ is
\begin{equation}
    V_{\sigma}^{(\mathrm{int})}(\mathbf r)
    =
    \frac{g}{(2\pi\sigma^2)^{D/2}}
    \exp\left(
        -\frac{\lVert\mathbf r\rVert^2}{2\sigma^2}
    \right),
\end{equation}
where $g$ is held fixed. The two protocols answer different questions and must
not be mixed in one convergence sequence. Each profile is periodized only
through a supercell large enough for image effects to be measured separately.

The ratio
\begin{equation}
    \rho=\frac{\sigma}{a}
\end{equation}
is one control, not a complete continuum criterion. The study must also record
the defect strength relative to the relevant bandwidth or gap, the bound-state
or envelope length relative to $a$ and the supercell, the retained-band and
valley separations, and the lattice and continuum discretizations. Frozen
sequences should range from atomically sharp profiles $\rho=O(1)$ to slowly
varying profiles $\rho\gg1$ under both normalization protocols.

For each $\rho$, compare:
\begin{enumerate}
    \item the exact represented lattice defect;
    \item a local lattice model;
    \item a finite-range nonlocal lattice model;
    \item a single-band effective-mass model where admissible; and
    \item multiband or multivalley continuum classes when the retained parent
          requires them.
\end{enumerate}
The study should measure operator residuals, band or valley mixing, bound-state
energies, subspace fidelity, exterior residuals, and cross couplings. Expected
improvement with increasing $\rho$ is a hypothesis to test, not a result to
assert in advance.

A proposed crossover radius or smoothness threshold is meaningful only when
both exterior and cross-coupling criteria pass and remain stable under lattice,
continuum-discretization, and supercell refinement. If no tested scale passes,
the correct outcome is no established continuum crossover over the tested
domain.

\section{Operator error versus bound-state error}

For a donor-like state below a declared conduction threshold, one possible
binding convention is
\begin{equation}
    E_b
    =
    E_{\mathrm{edge}}-E_{\mathrm{defect}}>0.
\end{equation}
An acceptor-like branch requires its own signed convention relative to the
aligned valence threshold. The parent and model must use the same declared
band-edge or continuum reference. At finite supercell size, both the defect
level and the relevant edge may depend on supercell momentum, so the comparison
must retain that dependence or state the rule used to extract an isolated-limit
quantity.

For every planted or fitted model, retain separately
\begin{align}
    \varepsilon_H
    &=
    \left\lVert
        \Delta\mathbf H_{
        \mathrm{parent}}
        -
        \Delta\mathbf H_{
        \mathrm{model}}
    \right\rVert_{
    \mathrm F},
    \\
    \Delta E_b
    &=
    E_{b,\mathrm{model}}
    -
    E_{b,\mathrm{parent}},
    \\
    F
    &=
    \left|
        \left\langle
            \psi_{b,\mathrm{parent}}
            \mid
            \psi_{b,\mathrm{model}}
        \right\rangle
    \right|^2
\end{align}
when individual normalized bound states are isolated, aligned, and present in
both models. Before matching states, report the number of in-gap or bound states
under the declared threshold convention. If the parent and model have unequal
counts, the mismatch is itself a failed diagnostic and no arbitrary one-to-one
fidelity should be forced. Degenerate or nearly degenerate clusters of equal
rank require a projector or principal-angle fidelity instead of an arbitrary
state overlap; unequal-rank clusters must retain both the rank mismatch and an
appropriately defined subspace discrepancy.

The one-dimensional synthetic exercise in
Appendix~\ref{app:impurity-model-classes} realizes both contrasting cases:
accurate binding with a large excited-sector operator residual, and a small
global residual with a significant bound-state change. Those observations
verify the metric separation for that constructed finite model only;
corresponding two-dimensional and material behavior remains untested.

\section{Null, covariance, and stopping conditions}

Every stage should include three outcomes:
\begin{description}
    \item[Null control.] With no planted defect, the aligned difference should
      vanish within the established representation and alignment errors.
    \item[Covariance control.] Under an admissible coordinate change, the
      abstract impurity operator should be unchanged and its matrix should
      transform by the declared conjugation rule.
    \item[Stopping control.] Incompatible rank, spin, geometry, subspace, or
      energy-reference data should prevent interpretation rather than produce a
      nominal residual.
\end{description}
These controls should be evaluated before any fitted model is called an
impurity potential or effective-mass representation.

\section{Required warmup order}

The proposed sequence is:
\begin{enumerate}
    \item verify the spinless-to-spinful embedding and Frobenius normalization;
    \item verify primitive-to-pristine-supercell folding;
    \item pass the zero-defect null extraction;
    \item recover known one-dimensional planted defects under gauge and energy
          scrambling;
    \item establish one-dimensional supercell-size convergence;
    \item recover the declared one-dimensional model-class hierarchy;
    \item scan one-dimensional defect smoothness toward a continuum regime;
    \item reproduce the same controls in the separable two-dimensional parent;
    \item introduce anisotropy and nonseparable two-dimensional coupling;
    \item compare operator, bound-state, and subspace metrics; and
    \item only then transfer the verified procedures to the matched
          pristine--doped silicon program.
\end{enumerate}
The two-dimensional extraction should not be interpreted until the
one-dimensional procedure recovers a known planted defect under alignment
scrambling and supercell refinement. Likewise, first-principles execution would
not be authorized or scientifically validated merely because the warmups pass.
As recorded in Appendix~\ref{app:impurity-model-classes}, the spin-space and
one-dimensional prerequisites have human-accepted calculated synthetic
evidence, including blind-alignment, independent-route, finite-rank-oracle, and
continuum-refinement follow-ons. The periodic two-dimensional folding parent is
also accepted. Under the version-1 defect design, Stage~A now supplies
human-accepted calculated and independently reproduced two-dimensional folding,
seam, null-extraction, and incompatibility controls. Nonzero two-dimensional defects in Stages~B--E,
new spin variants, and material transfer remain proposed and unauthorized.

\section{Remaining warmup plot placeholder}

The remaining placeholder is two-dimensional covariance: display extracted
defect blocks for symmetry-related placements after alignment, together with
the two-dimensional folding, null, supercell-shape, anisotropy, and metric
controls. The completed one-dimensional Figure~\ref{fig:impurity-defect-1d-benchmark}
replaces the previous one-dimensional placeholders. Every future plot should
identify the parent model, supercell, boundary condition, retained space, spin
convention, planted defect, alignment map, energy reference, model class,
normalization, and evidentiary status. A plot placeholder is not a calculated
result.

\section{Research questions}

The resulting program is organized by four questions.

First, which retained subspace and localized representation preserve the
declared pristine-silicon parent observables while remaining stable under
changes in numerical resolution and gauge?

Second, what is the least complex constrained lattice-model class that
satisfies both spectral and operator-level tolerances for the pristine host?

Third, which unitary identifications make pristine and doped retained operators
comparable, and how sensitive is the aligned impurity operator to admissible
choices of gauge, geometry, and energy reference?

Fourth, at what spatial scale, and within which anisotropic or multivalley model
class, can the aligned phosphorus or boron impurity operator be replaced by an
effective-mass description without losing the declared bound-state and operator
observables?

These questions place the continuum approximation at the end of the reduction
chain rather than assuming it at the outset. The aligned impurity operator is
the object extracted from first principles. A local screened potential, a
nonlocal lattice operator, and a multivalley effective-mass model are competing
admissible representations whose adequacy must be established by comparison.
